\section{双代数的左缩并和右缩并}

记号约定:$\left(E,c\right)$是分次双代数(相对的,反对称分次)$A$-双代数.

\begin{proposition}{}{}
	\begin{enumerate}
		\item 当$E$是分次双代数时,对每个$E$的每个$1$次齐次元$x$,映射$\iota_{x},\iota_{x}^{\prime}$是$\underline{E}^{\vee}$上的导子.
		\item 当$E$是反对称分次双代数时,对每个$E$的每个$1$次齐次元$x$,映射$\iota_{x},\iota_{x}^{\prime}$是$\underline{E}^{\vee}$上的反导子.
	\end{enumerate}
\end{proposition}

\begin{remark}{}{}
	设$N,N^{\prime},N^{\prime\prime}$是$\N$型分次$A$-模,$m\in\Bil_{A}\left(N,N^{\prime};N^{\prime\prime}\right)$.对每个$u\in\underline{\Hom}_{A}\left(E,N\right),v\in\underline{\Hom}_{A}\left(E,N^{\prime}\right)$,定义
	\begin{align*}
		u\cdot v\coloneq m\circ\left(u\otimes v\right)\circ c.
	\end{align*}
	设$x$是$E$中的$1$次元素,$u$是$\underline{\Hom}_{A}\left(E,N\right)$的$p$次齐次元素.
	\begin{enumerate}
		\item 若$E$为分次双代数,则$\iota_{x}\left(u\cdot v\right)=\iota_{x}\left(u\right)\cdot v+u\cdot\iota_{x}\left(v\right)$.
		\item 若$E$为反对称分次双代数,则$\iota_{x}\left(u\cdot v\right)=\iota_{x}\left(u\right)\cdot v+\left(-1\right)^{p}u\cdot\iota_{x}\left(v\right)$.
	\end{enumerate}
\end{remark}

\begin{remark}{}{}
	设$x\in E$,$I$是有限集且$c\left(x\right)=\sum\limits_{j\in I}x_{j}^{\prime}\otimes x_{j}^{\prime\prime}$,则对每个$u,v\in\underline{E}^{\vee}$有$\iota_{x}\left(uv\right)=\sum\limits_{j\in I}\iota_{x_{j}^{\prime}}\left(u\right)\cdot\iota_{x_{j}^{\prime\prime}}\left(v\right)$.特别地,当$E$是分次双代数,$c\left(x\right)=x\otimes 1+1\otimes x$时,$\iota_{x}$是导子.
\end{remark}






















